\begin{thebibliography}{99}
\bibitem{BPR} S. Basu, R. Pollack and M.-F. Roy,
\emph{Algorithms in Real Algebraic Geometry}, 2nd ed., Springer, 2006.
\bibitem{BB} I. Berkes and B. Borda,
Random walks on the circle and Diophantine approximation,
\emph{J. London Math. Soc.} \textbf{108} (2023), 409--440.
\bibitem{Cassels} J. W. S. Cassels,
\emph{An Introduction to Diophantine Approximation}, Cambridge Tracts in Mathematics and Mathematical Physics, vol. 45, Cambridge University Press, 1957.
\bibitem{DGLPS} J. Dick, T. Goda, G. Larcher, F. Pillichshammer and K. Suzuki,
On the quasi-uniformity properties of quasi-Monte Carlo point sets and sequences---Part I: Lattices and Kronecker sequences,
arXiv:2502.06202v2, version of 13 February 2026, Section 3.3.
\bibitem{GL} S. Graf and H. Luschgy,
\emph{Foundations of Quantization for Probability Distributions}, Lecture Notes in Mathematics, vol. 1730, Springer, 2000.
\bibitem{LPV} C. H. Lee, E. Pyne and S. Vadhan,
Fourier growth of regular branching programs,
\emph{APPROX/RANDOM 2022}, LIPIcs \textbf{245}, Art. 2, 2:1--2:21.
\bibitem{GTF35} Q. Qi,
\emph{General Theta Foundations I: Sharp Width of Hidden Rotation Experiments}, manuscript, revision 35, 2026; repository publication \texttt{aabc0731a0c9}.
\bibitem{GTF42} Q. Qi,
\emph{General Theta Foundations I: Alphabet Dynamics and Calibrated Numerical Memory}, manuscript, revision 42, 2026; repository publication \texttt{b5f9874903a4}.
\bibitem{RSV} O. Reingold, T. Steinke and S. Vadhan,
Pseudorandomness for regular branching programs via Fourier analysis,
\emph{RANDOM 2013}, LNCS \textbf{8096}, 655--670.
\bibitem{SVW} T. Steinke, S. Vadhan and A. Wan,
Pseudorandomness and Fourier growth bounds for width 3 branching programs,
\emph{Theory Comput.} \textbf{13} (2017), Art. 12, 1--50.
\bibitem{Vidyasagar} M. Vidyasagar,
\emph{Hidden Markov Processes: Theory and Applications to Biology}, Princeton University Press, 2014.

\bibitem{BF04} L. Benvenuti and L. Farina,
A tutorial on the positive realization problem,
\emph{IEEE Trans. Automat. Control} \textbf{49} (2004), 651--664,
Theorems 2--3 and Example 4.
\bibitem{Ben22} L. Benvenuti,
Minimal positive realizations: A survey,
\emph{Automatica} \textbf{143} (2022), Art. 110422.
\bibitem{BPP15} B. Balle, P. Panangaden and D. Precup,
A canonical form for weighted automata and applications to approximate minimization,
arXiv:1501.06841v3, 2015, Theorems 2 and 11.
\bibitem{FGS10} L. Finesso, A. Grassi and P. Spreij,
Two-step nonnegative matrix factorization algorithm for the approximate realization of hidden Markov models,
arXiv:1007.3435v1, 2010, Sections III--IV.
\bibitem{Ohta20} Y. Ohta,
On the realization of hidden Markov models and tensor decomposition,
arXiv:2008.11487v1, 2020, Theorems 6--9 and Remark 10.
\bibitem{BL06} Y. Bugeaud and M. Laurent,
Exponents of Diophantine approximation,
arXiv:math/0611354v1, 2006, Definition 1 and Section 2.
\bibitem{GTF33} Q. Qi,
\emph{General Theta Foundations I: Compatible Polyhedral Lifts and Quantitative Causal Width},
manuscript, revision 33, 2026; repository publication \texttt{aab961081253}.
\bibitem{GTF43} Q. Qi,
\emph{General Theta Foundations I: Word Profiles and Arithmetic Fluctuations of Numerical Memory},
manuscript, revision 43, 2026; repository publication \texttt{153830f5dc8d}.
\bibitem{Tzeng92} W.-G. Tzeng,
A polynomial-time algorithm for the equivalence of probabilistic automata,
\emph{SIAM J. Comput.} \textbf{21} (1992), 216--227,
doi:10.1137/0221017.
\bibitem{Shitov18} Y. Shitov,
A universality theorem for nonnegative matrix factorizations,
arXiv:1606.09068v2, version of 5 April 2018, Theorem 2 and Corollary 3.
\end{thebibliography}
